\scanpage{337}
\mspage{6}

\diagram{\begin{tikzcd}[ampersand replacement=\&,row sep=tiny]
 \& b\\
 b\times b\arrow[ur,"\mathrm{pr}_1"]\arrow[dr,"\mathrm{pr}_2"'] \&\\
 \& b
\end{tikzcd}}

dans $U'$, et les \unsure{hyp.|hypp.} \unsure{a)|b)}, b) de la page 1 assurent que la somme amalgamée \uline{dans} \unsure{de|des|du} ce diagramme \deleted{\unsure{du|de}}dans $U'^\wedge$ est $e_{U'}$, i.e. que

\[
 p_!(e_U)=e_{U'}.
\]

Il s'ensuit \deleted{\unsure{que}} ds la somme amalgamée $X'$ envisagée ds $(\alpha)$, l'unique élément $t'$ de $T'$ est objet final de $X'$. On a vu que $U'$ \deleted{\unsure{ne}} est 1-connexe, mais non 2-connexe si le groupe abélien enveloppant de $G$ (i.e. $G/$\inserted{groupe des} commutateurs dans le cas où $G$ est un groupe) n'est pas réduit à 1. L'exemple le plus économique s'obtient en prenant $G=\{\pm1\}$ le groupe à deux éléments. C'est donc un cas où dans \unsure{$(\alpha)$|cas} on a

\[
 \underset{\substack{\text{(et même}\\\text{asphérique)}}}{i\in W},\qquad \text{mais }i'\notin W,
\]

et qui a l'avantage sur l'exemple plus sophistiqué de la page \unsure{153|W3} (prop. \unsure{10|W}) de ne mettre en jeu que des catégories finies - et même fort petites !

Procédons maintenant comme dans la prop. 11, en désignant par $U''$ la catégorie déduite de $U'$ en lui ajoutant un objet

% archive sheet number: 337
